\section{Open questions}\label{sec:openq}

The following five questions are the highest-leverage unresolved items surfaced
by re-reading the paper against the analytic reproduction. Each is scoped so
that a targeted follow-up (either a modest MC job or a literature/data
comparison) could close it within one focused work session.

\begin{enumerate}
  \item \textbf{Plasmid-to-chromatin translation.} Does the FLASH sparing
        signature at the DSB level in naked pUC19 plasmid translate to
        eukaryotic chromatin, where histone shielding, higher local DNA
        density, and altered endogenous scavenger microenvironment
        ($\sigma\!\approx\!3$--$7\!\times\!10^{8}\ \mathrm{s^{-1}}$) materially
        change the intertrack $\cdot$OH kinetics? By the paper's own Fig.~4,
        the endogenous cellular $\sigma$ regime is precisely where the
        analytic model predicts \emph{no} UHDR sparing --- yet cell/animal
        FLASH experiments show sparing. Something is missing in either the
        plasmid model or in the mapping between physical DSB and biological
        endpoint.

  \item \textbf{Dose-rate crossover threshold.} What is the actual dose-rate
        threshold at which behavior transitions from CONV to UHDR, and is it
        a sharp threshold or a smooth crossover? The paper compares two dose
        rates eight orders of magnitude apart plus one bonus point; the
        interesting $10^{3}$--$10^{6}$~Gy/s decade (where most real FLASH
        beams operate) is not sampled.

  \item \textbf{In-vivo FLASH biomarker correlation.} Does the physical DSB
        yield signature predicted here correlate with any measurable in-vivo
        FLASH sparing biomarker ($\gamma$H2AX foci resolution, 53BP1 clearance,
        chromosome aberrations at G2), or is initial DSB yield simply the
        wrong endpoint for the FLASH effect? Published $\gamma$H2AX foci
        experiments at UHDR find minimal reduction in initial DSB, which
        would decouple this paper's physical prediction from clinical FLASH
        outcomes.

  \item \textbf{Cross-code TOPAS-nBio vs Geant4-DNA vs PARTRAC.} How does
        this specific TOPAS-nBio implementation (TsEmDNAPhysics + ELSEPA +
        Meesungnoen thermalization) compare quantitatively to Geant4-DNA
        standalone or PARTRAC on the identical plasmid geometry and pulse
        parameters? A cross-code CoV $>15\%$ on the DSB(UHDR)/DSB(CONV) ratio
        would call the specific 73.5\% number into question even while
        preserving the qualitative sparing prediction.

  \item \textbf{Mechanism decomposition: $\cdot$OH recombination vs oxygen
        chemistry.} Is UHDR DSB reduction driven primarily by $\cdot$OH
        intertrack recombination (paper's stated mechanism, R38*/R39*) or by
        transient oxygen depletion / fixation kinetics (Wilson-Alaghband,
        R40*/R41*)? The 20$\times$ difference between Model 1 (73.5\%) and
        Model 2 (3.5\%) reductions implies oxygen chemistry contributes
        substantial variance that the current analytic decomposition cannot
        cleanly separate.
\end{enumerate}

Machine-readable versions with citations and concrete next steps are in
\texttt{report/open\_questions.json}.
